% !TEX root = ../main.tex
\begin{abstract}
  The Bulk Synchronous Parallel (BSP) paradigm powers our most critical computational tools, from
  large-scale scientific simulations to distributed machine learning.
  As these applications scale, their efficiency has datacenter-level impacts on both cost and
  time-to-discovery.
  Yet the tightly synchronized nature of BSP creates two distinct challenges: on the data path,
  bursts of checkpointed state strain storage and analytics infrastructure, while on the compute
  path, global synchronization amplifies sensitivity to execution variability.
  While traditionally addressed by separate tools, this thesis demonstrates that both problems can be
  unified by treating them as challenges of organizing and processing streaming data for insight.

  We first address a critical data path challenge: enabling efficient ingestion for range queries
  over scientific data.
  Range queries are essential for scientific discovery, but require ordering via expensive
  postprocessing sorts that require multiple passes over on-disk data.
  An ideal solution would support single-pass, streaming ingestion that reorganizes data on the fly,
  even in the presence of highly dynamic and skewed key distributions.
  CARP provides this capability by reconstructing these distributions at runtime from rank-local
  statistics, and using it to inform in-situ reordering as data is written.
  % CARP provides this capability by adaptively tracking the evolving data distribution and performing
  % in-situ reordering as data is written.
  This allows CARP to approximate the query performance of a fully sorted dataset with no ingestion
  overhead, achieving up to $5\times$ faster end-to-end ingestion performance than prior approaches.
  % and up to $100\times$ lower query latency compared to other in-situ techniques.

  We then turn to the compute path, where we investigate low cluster efficiency in adaptive mesh
  codes.
  While often attributed to load imbalance, its underlying causes have not been systematically
  explored.
  Our analysis revealed that genuine load imbalance is frequently confounded by factors such as
  hardware faults, system noise, and unpredictable interactions across the system stack, all of which
  are amplified by global synchronization.
  To disentangle these effects, we developed custom, fine-grained telemetry and query-driven
  analytics pipelines.
  This process led to CPLX, a tunable placement policy that improves runtime by up to $21.6\%$.
  More importantly, it showed that the practical realization of compute efficiency depends as much on
  real-time visibility into runtime anomalies as on the placement algorithms themselves.
  % More importantly, it highlighted that More importantly, it highlighted the need for real-time,
  % queryable feedback to rapidly identify and address performance bottlenecks. This diagnostic
  % process led to CPLX, a tunable placement policy that improves runtime by up to $21.6\%$. More
  % importantly, it highlighted the need for real-time, queryable feedback to rapidly identify and
  % address performance bottlenecks.

  % Our experiences revealed that large-scale cluster efficiency is fundamentally bottlenecked by the
  % difficulty of reconstructing a coherent view from system-wide state that remains fragmented across
  % nodes and software layers. Existing methods for constructing global views---whether mesh
  % visualizations, performance traces, or data distributions---are brittle, slow, and specialized for
  % narrow domains. Our experience with custom telemetry pipelines demonstrated that relational
  % analytics could enable a unified approach to observability through programmable, query-driven
  % dataflows. ORCA realizes this vision, providing a streaming analytics framework that translates
  % observability requirements into optimized data movements between application and analysis tiers.
  % This unified substrate offers a programmable foundation to address longstanding efficiency
  % challenges across both data and compute paths in large-scale parallel systems.

  % We realized that both the CARP and AMR efforts ultimately relied on building specialized systems
  % to support a common capability: reconstructing distributed cluster views in real time. Modern
  % observability stacks—especially those centered on tracing—log vast volumes of event-level data for
  % offline analysis, but computing even basic cluster-wide summaries from such traces is
  % prohibitively expensive. Similar patterns exist in distributed machine learning workflows, where
  % fixed-function instrumentation and analysis of persisted traces are used for tasks ranging from
  % validating application correctness (e.g., detecting divergence or NaNs) to monitoring performance.
  % We propose a shift from this model: instead of recording everything upfront, treat execution as a
  % fragmented system image, where fragments are queried and composed on demand to reconstruct only
  % the views needed. ORCA realizes this approach through a tree-based overlay that leverages
  % declarative interfaces for programmable, on-demand aggregation of telemetry. It unifies
  % observability and control across both data and compute paths, enabling real-time introspection and
  % adaptation in large-scale parallel systems.

  While CARP and AMR targeted different problems, both revealed the same underlying constraint:
  meaningful efficiency gains were bottlenecked on the ability to rapidly reconstruct distributed
  application state.
  Today, BSP workloads across scientific computing and machine learning training rely on massive
  traces to identify correctness and performance bottlenecks.
  Tracing incurs significant runtime overheads, requires expensive postprocessing, and remains
  prohibitive at scale.
  We propose a fundamental shift: instead of proactive collection of all event logs, treat them as
  fragments of the global view, and use declarative dataflows for programmable, on-demand collection
  and aggregation of these fragments.
  ORCA enables these capabilities through an aggregation and control plane optimized for BSP
  telemetry structure, leveraging a tree-based overlay with SQL semantics and distributed query
  planning to minimize data movement.
  This replaces a fractured ecosystem of bespoke tools with a unified, query-driven
  framework for real-time observability across all layers and phases of distributed execution.

\end{abstract}
\newpage
